\scanpage{9}
\(\unsure{\varphi|\psi}\) soit sommable au voisinage de \(0\), on voit aussitôt qu'il faut et suffit que
\[
\int_1^\infty d\sigma\int_\sigma^\infty d\mu(s)<+\infty,
\qquad\text{i.e.}\quad\int_1^\infty(s-1)d\mu(s)<+\infty,
\]
\unsure{ce qui signifie} que
\[
\boxed{(6)\quad\int_1^\infty s\,d\mu(s)<\deleted{\unsure{\text{soit}}}+\infty}
\]
\deleted{\unsure{condition. Si} \(\mu=\mu_f\) (\(f\) fonction mesurable sur l'espace mesuré \(M\))}\inserted{Ceci est vérifié quand \(\unsure{\nu|f}\) est bornée, p.ex.} quand \(\nu=\mu_f\), \(f\) étant \unsure{comme dans les exemples, ou bornée}.

\deleted{Si alors \(\unsure{\mu_f}\) satisfait à \(\unsure{(6)}\)}
Si \(\varphi_f\) est sommable au voisinage de \(0\), on pose
\[
(7)\qquad\Phi_f(t)=\int_0^t\varphi_f(s)\,ds
\]
\(\Phi_f(t)\) est une fonction continue concave croissante de \(t\).
\inserted{\unsure{Considérons|En considérant} \(\log\unsure{\varphi_f(t)|\Phi_f(t)}\), c'est une fonction \deleted{décroissante}\inserted{\unsure{concave}}, à valeurs dans \([-\infty,+\infty[\).}
Enfin, supposons \(\log\varphi_f(t)\) sommable au voisinage de \(+0\) (\unsure{a fortiori} \(\varphi_f\) le soit),\inserted{i.e. \(\log f\) sommable dans \(f^{-1}([1,+\infty[)\)} \deleted{\(\log\varphi_f(t)\) \unsure{est}} Nous posons alors
\[
(8)\qquad\Psi_f(t)\unsure{=}\int_0^t\log\varphi_f(\unsure{t|s})
\]
Si \(\varphi_f\) ne s'annule pas \deleted{\unsure{limite à} \(\unsure{+\infty}\) \illegible{2}} (i.e. si \(f^{-1}(\unsure{\complement\{0\}})\) n'est pas \uline{intégrable}) \(\Psi_f(t)\) est une fonction partout finie. Sinon, si \(t_0=|f^{-1}(\complement\{0\})|\) est le \(t\) à partir duquel \(\varphi_f\) est nulle, c'est le point à partir duquel \(\Psi_f(t)\) est égale à \(-\infty\).

Enfin, on pose
\[
(9)\qquad\Delta_f(t)=\exp\Psi_f(t)
\]
Sous les conditions envisagées sur \(f\), \(\Delta_f(t)\) est une fonction continue \deleted{positive}\inserted{\unsure{convexe|concave}} de \(t\); elle ne peut s'annuler que si \(t_0=|f^{-1}(\complement\{0\})|\) est fini, alors \(\Delta_f(t)\) s'annule constamment à partir de \(t_0\).

Les formules (5) donnent des inégalités correspondantes pour \(\Phi_f,\ \Psi_f,\ \Delta_f\). La plus importante est
\[
(10)\qquad\Psi_{fg}\leq\Psi_f+\Psi_g\qquad\text{i.e.}
\]
\[
(11)\qquad\Delta_{fg}\leq\Delta_f\Delta_g
\]

\mspage{10}
